Ci sono molte persone che desidero ringraziare alla fine di questo percorso, tutte le quali sono state fondamentali
per arrivare, oggi, a questo traguardo.

Innanzitutto voglio ringraziare il mio relatore, il professor Daniele Carnevale, il quale mi ha concesso l'opportunità di lavorare
a questo interessantissimo progetto di tesi e ha creduto nelle mie capacità più di quanto facessi io stessa.

Ringrazio anche Andrea Locatelli, che, ancora prima di sapere di essere il mio correlatore ufficiale,
ha sopportato le mie infinite domande e mi ha fornito un
aiuto fondamentale affinché questo progetto funzionasse e la mia tesi vedesse, finalmente, la luce. 

Un ringraziamento particolare va a Edoardo, che, come un mentore, mi ha ispirata e guidata attraverso le insidie di 
questo percorso universitario, invitandomi sempre a non scoraggiarmi. Sei stato un ottimo amico, anche nella stoica 
sopportazione delle mie lamentale.

Ringrazio Sara, Silvia e Martina, che hanno accompagnato il mio percorso sin dalla triennale e nelle quali ho trovato delle
amiche oltre che delle colleghe; grazie per le risate, gli attacchi di ansia poi sedati e tutti gli appunti che avete preso per 
me quando non ero a lezione. Senza le vostre spiegazioni e il supporto reciproco, avrei ancora esami da sostenere e il tempo 
passato in università non sarebbe stato altrettanto divertente.

Ringrazio Ilaria e Alessandra, con le quali ho condiviso tutti i giorni della magistrale. In questi due anni di lezioni ho passato più 
tempo insieme a voi che con la mia famiglia, ma sicuramente siete stata una piacevole compagnia. Studiare e lavorare insieme ad ogni progetto
mi ha permesso di apprendere tanto e ha portato una ventata di leggerezza nelle grigie giornate piene di calcoli e righe di codice da scrivere.

Ringrazio tutte le persone che ho conosciuto al Campus, che hanno accompagnato la mia vita, universitaria e non, in questi anni e che hanno 
rappresentato una seconda famiglia in un posto nuovo che ho imparato a considerare una seconda casa.

In particolare, voglio ringraziare Ugo, che ho conosciuto ben prima di arrivare qui e che è diventato uno dei miei amici più cari.
Ti ringrazio di tutte le chiacchierate fino a tarda notte, anche quando volevo solo andare a dormire e ti dovevo praticamente cacciare da casa mia,
di tutti i momenti condivisi, belli e meno belli, di questi anni, delle vacanze trascorse insieme e di tutto il supporto che mi hai sempre dato.

Grazie a Caterina, anche lei mia amica da quasi sei anni. Abbiamo vissuto fianco a fianco tutto il percorso universitario, ci siamo viste nelle condizioni peggiori, 
con indosso maglioni improbabili, e abbiamo imparato a capirci con uno sguardo. Anche se non vivremo più insieme, so che abbiamo condiviso un legame speciale, che, spero, continuerà a crescere anche
nel futuro.

Grazie a Francesca, che è un'amica dolcissima e sempre pronta a consolarmi o a gioire insieme a me. Grazie di ascoltare sempre
ciò che ho da dire senza giudizio e di essere stata una presenza discreta ma tangibile in alcuni momenti difficili.

Ringrazio Giovanni, a cui dimostro il mio affetto con le battute antipatiche, per essere stato un amico disponibile e solare, grazie 
al quale non ci si annoia mai. Grazie anche di tutte le cose buone che mi hai fatto mangiare, che mi hanno ricordato casa.

Un riconoscimento va a Grazia, che allieta le mie giornate con la sua dolcezza e il suo umorismo. Grazie per il nostro club del libro, 
che mi ha permesso di scoprire tante letture nuove, e di avermi insegnato l'arte di procurarmi in maniere alternative prodotti letterari e audio-visivi
protetti da copyright.

Ringrazio tutto il coro della Cappella Universitaria San Tommaso, di cui molte delle persone già citate fanno, o hanno fatto, parte,
perché mi accompagna sin dall'inizio di quest'avventura.
Devo molto a tutti loro, poiché mi hanno permesso di coltivare una mia passione e sono diventati miei amici, cantando e condividendo le gioie 
e i dolori dovuti a impegni ed esami universitari, accompagnandomi nel percorso di crescita che mi ha condotta fin qui.

Un ringraziamento speciale va alla mia famiglia, che in questi anni non è stata mai troppo lontana. Ci sarebbe un'infinità di cose
da dire, ma devo condensare tutto l'affetto e la gratitudine in poche parole.

Grazie a mamma e papà per tutto l'amore che mi avete sempre dato, per aver creduto nelle mie capacità e 
per la vicinanza costante anche a chilometri di distanza. Rappresentate sempre i miei punti di riferimento,
a cui so di poter tornare in qualsiasi momento, pronta ad essere accolta a braccia aperte.

Grazie a Beatrice, che è sempre stata una perfetta sorella maggiore, soprattutto quando mi sono trasferita a Roma e 
ha rappresentato il ponte tra la mia vecchia vita e la nuova. Grazie di tutto quello che abbiamo condiviso: i libri,
i film, le esperienze, le discussioni, il coro stesso, a cui non avrei mai preso parte senza di te. 

Grazie a Elisabetta, con la quale non ci siamo sempre capite, ma a cui voglio comunque un mondo di bene. Abbiamo trovato 
il nostro equilibrio, anche nella convivenza, e sei stata un supporto fondamentale sia per smorzare la mia tendenza all'ansia 
che per affrontare alti e bassi della vita.

Grazie a Pietro, il mio, non più così piccolo, fratellino, con il quale riesco a fare sempre discorsi filosofici che elevano 
il mio intelletto e mi consentono di usare il mio cervello in maniera diversa dal ragionamento matematico a cui ormai è abituato.
Grazie anche di avermi fornito uno strumento prezioso per ravvivare la mia spirtualità, che va di pari passo con il mio benessere psico-fisico.

Ringrazio tutte le persone che ho incontrato durante il periodo universitario, professori, colleghi, amici, persone che hanno fatto parte della mia vita
solo per un periodo. Questo percorso mi ha cambiata, ho avuto il privilegio di apprendere un'infinità di cose e di fare una 
moltitudine di esperienze. Ho imparato a credere di più nelle mie capacità, soprattutto perché tante persone ci hanno creduto al 
posto mio quando io mi abbattevo, e ho capito che bisogna dare il giusto peso alle cose: che fare delle domande non significare sembrare
stupidi e un esame andato male non è la fine del mondo.

Sono grata di essere giunta al termine di questo viaggio, impegnativo ed entusiasmante, e di indossare questo alloro come simbolo del coronamento
di un sogno a lungo agognato.


